\section{Further questions}\label{sfh:sec:questions}
The results isolate several extensions with distinct analytic obstacles.

\paragraph{Effective constants and certified integer inversion.}
The two-ceiling theorem becomes a usable certified large-range threshold procedure once the full-arc and local Taylor constants are explicit beyond a concrete saddle. It would be useful to combine interval evaluation of $F_J$ with a provable tail integral bound and exact adjacent recurrence checks. Finite numerical agreement does not provide those constants. An implementation should also control cancellation when evaluating the small difference $V=C-n^2/b$.

\paragraph{Higher conditional likelihood corrections.}
The local Fourier proof can retain higher weighted degrees, yielding polynomial corrections in the centered removed mass. The main additional issue is to organize these corrections so that their signed $Q$-mass cancels order by order and their global remainder stays summable under a growing-parameter compound Poisson law. Such an expansion could refine the sharp total-variation equivalent beyond its first constant. It would also be natural to turn a corrected signed approximation into a genuine probability law without losing the claimed order.

\paragraph{An optimal hybrid error.}
The bound $O(r^3/\sqrt n)$ in Theorem~\ref{sfh:thm:hybrid} is driven by the exceptional-count contribution to the conditional component mean. One may ask whether explicitly centering $K$ by that observed contribution gives a smaller uniform mixed-test error. A result in a stronger metric would need to respect the lattice nature of $K$ rather than compare it in total variation directly to a continuous law.

\paragraph{Moderate and large edge deviations.}
The two Gaussian scales describe bounded normalized fluctuations. Tilting the defect marker may change the balance between the rare tetrahedral components and the ordinary pair-stars. A uniform marked saddle analysis with growing Fourier or exponential parameters could locate the crossover between Gaussian fluctuations, Poisson moderate deviations, and genuine large deviations. The compact-parameter proof here does not imply such a result.

\paragraph{Larger uniformity and structural obstructions.}
For $3$-uniform systems the component classification collapses to pair-stars and two finite exceptions. For higher uniformity, forbidden singleton intersections admit more complicated connected families; the simple polynomial-times-exponential connected EGF may disappear. Determining which restricted connected subclasses still admit a comparable explicit Poisson correction would connect enumeration more directly with the structural questions in the extremal literature.

\paragraph{Growing deletions.}
Theorem~\ref{sfh:thm:deletion} fixes the deleted polynomial independently of $n$. Allowing its degree or coefficients to grow would test the boundary of the small saddle displacement $p_1/b$. A uniform theorem would need both a quantitative nonnegative-coefficient condition and control of how the deleted mass changes every Fourier arc. It cannot be obtained merely by replacing the fixed degree in the displayed big-$O$ bound by a function of $n$.
\begin{writenote}[Added 7 October 2026, under Vladimir's standing rule of 4 October 2026]
These six directions are the source's own and stay open. The write adds from
the non-claims explicit constants for Theorem~\ref{sfh:thm:inverse} and
certified diagnostics. No claim of the source was found false, and the OEIS
entries contain no conjecture.
\end{writenote}
